% Appendices A and B of the thesis, parked on 8 October 2026 (removed from main.tex, not deleted).
%
% To add them back:
%   1. In main.tex, put  % Appendices A and B of the thesis, parked on 8 October 2026 (removed from main.tex, not deleted).
%
% To add them back:
%   1. In main.tex, put  % Appendices A and B of the thesis, parked on 8 October 2026 (removed from main.tex, not deleted).
%
% To add them back:
%   1. In main.tex, put  % Appendices A and B of the thesis, parked on 8 October 2026 (removed from main.tex, not deleted).
%
% To add them back:
%   1. In main.tex, put  \input{parked/appendices}  on the line before  \end{document}.
%   2. Move the two reference entries at the end of this file back into the References list of main.tex
%      (alphabetical order: after "Department of ..." entries / before "European Commission ..." etc.).
%   3. Restore the cross-references in the main text: search this file's labels app:uk2015 and app:earlier
%      (and sec:statistics, sec:baselines, sec:confirmatory, tab:baselines, tab:confirmatory,
%      tab:equivalence, tab:efficiency-tests, tab:uk-systems) and point the main text to them again
%      where it now says "the project repository" or names the results directly.
%
% Sources: results/uk2015.jsonl, results/tools/scenarios_uk2015.jsonl (Appendix A); results/runs.jsonl,
% results/baselines.jsonl, results/confirmation.jsonl, results/confirmation/confirmatory_tests.csv (Appendix B).

% ════════════════════════════════════════════════════════════════════════════
\appendix
\renewcommand{\thechapter}{\Alph{chapter}}

\chapter{Robustness check: UK prices of the REFIT period}
\label{app:uk2015}

% Data: data/scenarios_uk2015 (manifest), data/raw/prices_gb_entsoe.csv; results:
% results/uk2015.jsonl (configs/uk2015.json) and results/tools/scenarios_uk2015.jsonl (EMHASS).
The benchmarks of this thesis combine appliance cycles recorded in UK
homes in 2013--2015 with Finnish prices of 2022--2026
(Section~\ref{sec:dataset}). This appendix checks whether the results
also hold when the appliance data are priced in their own country and
period. It is a robustness check that was not pre-registered, and its
results are exploratory.

\section{Prices, tariff and scenarios}

The ENTSO-E Transparency Platform publishes day-ahead prices for the
Great Britain bidding zone only from 17 December 2014, in GBP per MWh
(ENTSO-E, 2026). This covers the last months of the REFIT recordings,
which ended in mid-2015 (Murray et al., 2017). The prices of 18 December
2014 to 31 July 2015, hourly, give 224 complete days on London time.
They were converted to euros at the mean of the European Central Bank's
daily reference rates for the same period, 0.73161 GBP per EUR
(European Central Bank, 2026), so that the agents' tools and the costs
use the same currency as in the other benchmarks. The saving in percent
does not depend on the currency.

The consumer price follows the form of the Finnish tariff
(Section~\ref{sec:dataset}): the spot price plus a constant charge per
kilowatt-hour, with VAT. The constant charge stands for all
non-wholesale costs of a UK household in 2015, that is network costs,
policy costs, the supplier's costs and margin, and the standing charge
spread over consumption. It was set so that the average consumer price
equals the average UK unit cost of standard electricity in 2015, 15.4
pence per kWh including 5\,\% VAT, which follows from DECC's average bill
of \pounds 584 for an annual consumption of 3\,800\,kWh (Department of
Energy \& Climate Change, 2016). With a mean spot price of 4.05 pence per
kWh in the window, the charge is 10.62 pence per kWh, or 0.1452 euros
per kWh. In 2014--2015 almost no UK household had a dynamic hourly
tariff, so the set describes a hypothetical household on such a tariff
rather than a real one.

The scenarios follow the factorial design of
Section~\ref{sec:factorial} with one change. The window contains no day
with a negative price, so the set has three price patterns, stable,
volatile and spike, each defined relative to the days of the window.
With two household loads, two peak limits and four days per pattern and
load, the set has 48 scenarios, 14 of which require coordination. The
cost of starting every appliance when it is ready has a median of
\EUR{0.63} per scenario. Every system of the confirmatory evaluation was
run on the set in the same way as on Q2: the Multi-Agent Pipeline and
the LLM without tools three times per scenario with GPT-4o-mini, and
EMHASS and the rule-based systems once.

\section{Results}

Table~\ref{tab:uk-systems} summarises the results. The ranking of the
systems is the same as on the Finnish benchmarks. The pipeline produced
a valid schedule in all 144 runs and saved 6.69\,\% of the cost without
scheduling, the same as the optimal schedule. EMHASS saved 6.45\,\% and
the smart timer 6.55\,\%. The cheapest-hours automations saved 2.6 to
3.6\,\%, and on the 14 scenarios that require coordination only 7\,\% of
their schedules were valid, against 100\,\% for the pipeline, EMHASS and
the smart timer. The LLM without tools saved 0.18\,\%, and a quarter of
its runs cost more than starting every appliance when ready. The
pipeline's mean latency was 18.5 seconds, and its API cost \$0.0025 per
recommendation.

\begin{table}[h!]
\caption{Robustness set: UK prices and tariff, December 2014 to July 2015 (48 scenarios; saving pooled over all runs).}
\label{tab:uk-systems}
\small
\begin{tabularx}{\textwidth}{>{\raggedright\arraybackslash}Xcccc}
\toprule
System & Valid & Saving & \makecell{Saving,\\strict limit} & \makecell{Saving,\\realistic limit} \\
\midrule
Multi-Agent Pipeline & 100\,\% & 6.69\,\% & 6.4\,\% & 6.9\,\% \\
Smart timer (own heuristic) & 100\,\% & 6.55\,\% & 6.2\,\% & 6.9\,\% \\
EMHASS, measured profiles & 100\,\% & 6.45\,\% & 6.0\,\% & 6.9\,\% \\
EMHASS, nominal power & 100\,\% & 6.06\,\% & 5.6\,\% & 6.5\,\% \\
Cheapest window & 70.8\,\% & 3.55\,\% & 1.5\,\% & 5.6\,\% \\
Cheapest-hours automation & 68.8\,\% & 2.62\,\% & 0.9\,\% & 4.3\,\% \\
LLM without tools & 95.1\,\% & 0.18\,\% & 0.3\,\% & 0.1\,\% \\
\midrule
Optimum (upper bound) & 100\,\% & 6.69\,\% & 6.4\,\% & 6.9\,\% \\
\bottomrule
\end{tabularx}
\end{table}

Paired per scenario, the pipeline saved more than the cheapest-hours
automation (by 3.3 percentage points), the cheapest window (2.1) and the
LLM without tools (6.4), with Wilcoxon $p < 0.001$ in each case, and it
saved the same as EMHASS within $\pm 1$ percentage point (mean difference
$+0.20$ points, 90\,\% confidence interval $[0.06, 0.34]$). These tests
mirror the registered ones but were not registered for this set, so they
are reported without adjustment.

The savings are about a quarter of those on the Finnish benchmarks.
Because the price patterns are defined relative to the days of the
window, even the volatile and spike days of 2015 are mild compared with
the Finnish ones: on these days the spot price varied within the day by
8 to 9 cents per kWh on average, against 16 to 20 cents on the
corresponding days of Q2. The window also has no negative prices. In addition,
the constant charge makes up most of the consumer price, so a given
shift in time changes a smaller share of the cost. The order of the
price patterns also differs from the Finnish sets: the stable days gave
the largest saving (9.8\,\% at the optimum) and the volatile days the
smallest (3.0\,\%). The reason was not analysed further.

\section{Conclusion}

Priced at the prices and tariff of their own country and period, the
REFIT appliance cycles give much smaller savings, but the comparison
between the systems does not change. The pipeline again reaches the
optimal saving with a valid schedule every time, matches EMHASS, and
saves more than the cheapest-hours automations, whose schedules again
fail where the household's appliances must be staggered. The main
results of this thesis therefore do not depend on combining UK
appliance data with Finnish prices, although the size of the saving
does depend on the price conditions.


% ════════════════════════════════════════════════════════════════════════════
\chapter{Evaluation of the earlier design}
\label{app:earlier}

This appendix reports the evaluation of the earlier design of the study,
summarised in Section~\ref{sec:earlier-design}. It compared the base
multi-agent pipeline, without retrieval or the optimiser tool, with a
single ReAct agent that had the first seven tools of
Table~\ref{tab:tools} and the same step budget, on the first benchmark of
48 scenarios in six types (Section~\ref{sec:dataset}) and on the
held-out Q set (Section~\ref{sec:resolution}). Its outcome was the
optimality gap, the cost of a schedule above the optimal cost in percent
of the optimal cost, with an infeasible schedule scored at the cost of
starting every appliance when ready. The research questions and
hypotheses of this design were replaced in the redesign, so the labels
SQ1, SQ2, H1 and H2 in Section~\ref{sec:confirmatory} refer to the first
pre-registration. The results remain valid for that design.

\section{Pipeline and single agent on the first benchmark}
\label{sec:statistics}

% Hand-typed from `hems analyze` (results/tables/equivalence.csv, efficiency_tests.csv), GPT-4o-mini only.
This section reports the tests of the earlier design for GPT-4o-mini on
the first benchmark, with the mean over three runs per scenario as the
paired observation ($N = 48$). These results are exploratory; the
pre-registered confirmatory results are reported in
Section~\ref{sec:confirmatory}. A run with the local model gpt-oss-20b
was stopped before completion and is not reported.

\paragraph{Schedule quality.} The multi-agent pipeline reached a mean
optimality gap of 1.25\,\% and the single agent 1.43\,\%. The Wilcoxon
signed-rank test finds no significant difference ($W = 449$,
$Z = -0.14$, $p = 0.89$, $r = 0.02$), and the median paired difference
is zero, because both designs reach the same cost in most scenarios. The
equivalence test shows that this is not merely a lack of evidence. The
mean paired difference is $-0.18$ percentage points, with a 90\,\%
confidence interval of $[-0.71, 0.35]$, which lies inside the primary
margin of $\pm 1$ percentage point (Table~\ref{tab:equivalence}). The two
designs are therefore equivalent in schedule quality within one
percentage point of the optimum. At the narrower margin of 0.5
percentage points the $t$ version of the test does not conclude
equivalence, because a few scenarios with large differences in either
direction widen the interval, although the Wilcoxon version does.

\begin{table}[h!]
\caption{Equivalence test (TOST) of the optimality gap, multi-agent vs.\ single-agent (GPT-4o-mini, $N = 48$).}
\label{tab:equivalence}
\small
\begin{tabular*}{\textwidth}{@{\extracolsep{\fill}}lccc}
\toprule
Margin & $p$, paired $t$ TOST & $p$, Wilcoxon TOST & Equivalent \\
\midrule
$\pm 0.5$ percentage points & 0.156 & $<$\,0.001 & No \\
$\pm 1$ percentage point (primary) & 0.006 & $<$\,0.001 & Yes \\
$\pm 2$ percentage points & $<$\,0.001 & $<$\,0.001 & Yes \\
\bottomrule
\end{tabular*}
\end{table}

\paragraph{Efficiency.} The two designs differ significantly in every
efficiency measure (Table~\ref{tab:efficiency-tests}). The pipeline used
34\,\% fewer input tokens, because each specialist works with a short
context and only its own tools, while the single agent carries the
entire conversation and all eight tool descriptions through every step.
It produced 83\,\% more output tokens and made 23\,\% more tool calls,
since four agents each write their own structured output and call their
tools separately. Input tokens dominate the price, so the API cost per run was
21\,\% lower for the pipeline, and lower in 35 of the 48 scenarios. The
sequential hand-offs made the pipeline 69\,\% slower, with a mean of
21.4 seconds against 12.7 seconds, and it was faster in only two
scenarios.

\begin{table}[h!]
\caption{Efficiency per run, multi-agent vs.\ single-agent (GPT-4o-mini; tokens and tool calls per run, API cost in USD, latency in seconds; each value is the mean over $N = 48$ scenarios of the per-scenario mean of three runs; Wilcoxon signed-rank test).}
\label{tab:efficiency-tests}
\small
\begin{tabular*}{\textwidth}{@{\extracolsep{\fill}}lcccccc}
\toprule
Measure & \makecell{Single\\agent} & \makecell{Multi-\\agent} & Difference & \makecell{Multi-agent\\lower} & $p$ & $r$ \\
\midrule
Input tokens & 21\,686 & 14\,292 & $-34$\,\% & 43/48 & $<$\,0.001 & 0.79 \\
Output tokens & 661 & 1\,208 & $+83$\,\% & 1/48 & $<$\,0.001 & 0.87 \\
Total tokens & 22\,347 & 15\,500 & $-31$\,\% & 40/48 & $<$\,0.001 & 0.75 \\
API cost (USD) & 0.0037 & 0.0029 & $-21$\,\% & 35/48 & $<$\,0.001 & 0.54 \\
Latency (s) & 12.7 & 21.4 & $+69$\,\% & 2/48 & $<$\,0.001 & 0.86 \\
Tool calls & 16.5 & 20.3 & $+23$\,\% & 1/48 & $<$\,0.001 & 0.81 \\
\bottomrule
\end{tabular*}
\end{table}

Taken together, decomposing the task into specialist agents did not
lower the cost of the schedules, but it reached the same quality with
about a third fewer tokens and a fifth lower inference cost, at the price
of a longer response time.

\section{Comparison with baselines and real tools}
\label{sec:baselines}

% Hand-typed from `hems analyze --runs results/runs.jsonl results/baselines.jsonl` (GPT-4o-mini).
Table~\ref{tab:baselines} compares the two agent designs with the smart
timer, the LLM without tools and the two real home energy management
tools. All systems that use a language model use GPT-4o-mini.

% EMHASS rows: Docker/REST service with default settings (tools/emhass_rest.py, systems emhass_rest*); cheapest hours;
% all in results/tools/scenarios.jsonl (secondary, added after the pre-registration). Savings captured = pooled.
\begin{table}[h!]
\caption{Agent designs, simpler baselines and real home energy management tools (GPT-4o-mini, main benchmark, 48 scenarios).}
\label{tab:baselines}
\small
\begin{tabular*}{\textwidth}{@{\extracolsep{\fill}}lcccccc}
\toprule
System & \makecell{Feasible\\runs} & \makecell{Mean\\gap} & \makecell{Median\\gap} & \makecell{Gap,\\type D} & \makecell{Savings\\captured} & \makecell{Time\\per run} \\
\midrule
Smart timer (own heuristic) & 48/48 & 0.23\,\% & 0.00\,\% & 1.2\,\% & 99\,\% & $<$\,0.01\,s \\
EMHASS, profiles & 48/48 & 0.32\,\% & 0.00\,\% & 0.1\,\% & 99\,\% & 1.2\,s \\
Multi-agent pipeline & 144/144 & 1.25\,\% & 0.01\,\% & 5.9\,\% & 93\,\% & 21.4\,s \\
Single agent & 143/144 & 1.43\,\% & 0.00\,\% & 6.5\,\% & 90\,\% & 12.7\,s \\
EMHASS, nominal power & 48/48 & 2.24\,\% & 0.45\,\% & 8.5\,\% & 87\,\% & 1.0\,s \\
Cheapest hours & 25/48 & 17.6\,\% & 4.2\,\% & 62.7\,\% & 25\,\% & $<$\,0.01\,s \\
LLM without tools & 126/144 & 23.1\,\% & 11.9\,\% & 66.0\,\% & 13\,\% & 2.0\,s \\
\bottomrule
\end{tabular*}
\end{table}

Without tools, GPT-4o-mini schedules poorly. Its mean optimality gap is
23.1\,\%, and 18 of its 144 schedules are infeasible: 14 exceed the peak
limit and 4 start a cycle outside its window. A further 22 of its
feasible schedules cost more than starting every appliance as soon as it
is ready, and over all its runs it captures only 13\,\% of the
achievable saving, against about 90\,\% for the agents. Both agent designs are significantly better (Wilcoxon
$p < 0.001$, $r = 0.87$ for each), with median paired differences of
10.8 percentage points for the single agent and 11.4 for the pipeline.
The same model that fails to schedule from a single prompt comes close to
the optimum when it can query prices, costs and loads through tools and
check its own work. Tool use, rather than the language model's own
calculation, is therefore what makes the agents viable.

The smart timer, however, outperforms both agent designs. It finds a
feasible schedule in every scenario, is optimal in 30 of the 48, and its
largest gap is 3.7\,\% (D04). Both agent designs have significantly
larger gaps than the timer (single agent $p < 0.001$, $r = 0.63$;
pipeline $p < 0.001$, $r = 0.68$). The difference is concentrated in
type D, where the timer's mean gap is 1.2\,\% against 5.9\,\% for the
pipeline and 6.5\,\% for the single agent. The median paired differences
are small (0.002 percentage points for the single agent and 0.017 for the
pipeline), because in most scenarios all three systems reach the optimum
or come very close to it. The timer needs no language model, runs in
about two milliseconds and has no inference cost. In this benchmark,
where the power profile and window of every appliance are known exactly
and the horizon is one day, placing the appliances one at a time is
almost always enough. The agents come within about one percentage point
of this rule but do not improve on it. What they add is the ability to
start from a request written in natural language rather than from
structured appliance data, and to explain the schedule they produce.

The two real home energy management tools, added after the
pre-registration, bracket the agents. Configured with the measured power
profiles, EMHASS comes close to the optimum: every schedule is feasible,
32 of the 48 are optimal, and the mean gap is 0.32\,\%. Its largest gap,
4.2\,\% in F02, comes from its default solver tolerance of 1\,\%
(the setting \texttt{lp\_solver\_mip\_rel\_gap} = 0.01; EMHASS
contributors, 2026), which
applies to the whole electricity bill including the base load and can
therefore amount to several percent of the appliance cost scored here.
With an exact tolerance the same optimization comes within 0.2\,\% of
the optimum in every scenario, which confirms that the MILP reference reflects what a
deployed open-source optimizer can achieve. The agents' mean gaps are
about one percentage point higher, but the difference from EMHASS as
deployed is not significant (Wilcoxon $p = 0.24$ for the single agent and
$p = 0.19$ for the pipeline). Configured with a constant nominal power,
EMHASS still keeps every schedule within the peak limit, because it
plans each cycle at its peak power, but it misjudges the cost of the real
cycles and reaches a mean gap of 2.2\,\%. Both agent designs are
significantly better than this configuration (Wilcoxon $p < 0.001$ for
each, $r = 0.67$ for the single agent and $0.64$ for the pipeline). The
cheapest-hours automation, under the assumption stated in
Section~\ref{sec:evaluation} that its trigger starts a cycle which then
runs to completion, exceeds the peak limit in 23 of the 48 scenarios,
because each appliance follows
its own cheap period and appliances whose windows share a price minimum
start together. Scored at the naive cost when infeasible, it reaches a
mean gap of 17.6\,\%, and both agent designs are significantly better
($p < 0.001$, $r = 0.82$ and $0.81$). The agents are therefore
outperformed by the smart timer and are not significantly different from
a well-configured optimizer as deployed, while they outperform a simply
configured optimizer and the per-device automation that is common among
households.

Two checks show that this ranking does not depend on how the costs are
scored. First, scoring an infeasible run at the cost of the naive
schedule does not drive the result. Restricted to feasible runs, the
mean gaps are 0.23\,\% for the smart timer, 0.32\,\% for EMHASS with
profiles, 1.25\,\% for the pipeline, 1.45\,\% for the single agent,
2.24\,\% for EMHASS with nominal power, 7.9\,\% for the cheapest-hours
automation and 23.0\,\% for the LLM without tools, the same order as in
Table~\ref{tab:baselines}, and the order on the held-out benchmark is
also unchanged. The cheapest-hours automation looks better on feasible
runs only because these are the days on which its appliances happen not
to collide. Second, the ranking is the same under the spot price alone
as under the consumer price, as explained in
Section~\ref{sec:evaluation}.

\section{Confirmatory run on the held-out Q set}
\label{sec:confirmatory}

% Hand-typed from `hems confirm` (results/confirmation/confirmatory_tests.csv) and
% `hems analyze --runs results/confirmation.jsonl` for the secondary results.
The hypotheses of the earlier SQ1 and SQ2 were fixed in the pre-registration, which
was committed to the project repository on 30 September 2026 together
with the code, prompts, tools and analysis script. The confirmatory run
followed on the held-out benchmark of 48 scenarios with fifteen-minute
prices (Section~\ref{sec:resolution}), on which no language model had
been run before. All systems that use a language model used GPT-4o-mini
with three runs per scenario, and the analysis was run exactly as
registered. The run cost USD~0.92 in API fees. Table~\ref{tab:confirmatory}
reports the six registered tests.

\begin{table}[h!]
\caption{Pre-registered confirmatory tests on the held-out fifteen-minute benchmark (GPT-4o-mini, $N = 48$; $p$-values Holm-adjusted within each sub-question).}
\label{tab:confirmatory}
\small
\begin{tabularx}{\textwidth}{l>{\raggedright\arraybackslash}Xcccc}
\toprule
& Hypothesis & Estimate & $p$ & $r$ & Supported \\
\midrule
H1 & Single agent gap $<$ LLM without tools & $-17.9$ pp & $<$\,0.001 & 0.86 & Yes \\
H1 & Pipeline gap $<$ LLM without tools & $-17.7$ pp & $<$\,0.001 & 0.86 & Yes \\
H2a & Pipeline equivalent to single agent ($\pm 1$ pp) & $+0.22$ pp & $<$\,0.001 & -- & Yes \\
H2b & Pipeline uses fewer input tokens & $-31$\,\% & $<$\,0.001 & 0.84 & Yes \\
H2b & Pipeline has a lower API cost per run & $-21$\,\% & $<$\,0.001 & 0.70 & Yes \\
H2c & Pipeline has a longer latency & $+52$\,\% & $<$\,0.001 & 0.80 & Yes \\
\bottomrule
\end{tabularx}
\end{table}

All six hypotheses are supported. For H1, both agent designs reached a
mean optimality gap of about one percent (0.86\,\% for the single agent
and 1.08\,\% for the pipeline), whereas the LLM without tools reached
18.8\,\%. Fifteen of its 144 schedules were infeasible, and over all its
runs it captured only 24\,\% of the achievable saving, against 92\,\% for
the single agent and 90\,\% for the pipeline. For H2a,
the mean paired difference in gap was $+0.22$ percentage points with a
90\,\% confidence interval of $[-0.15, 0.59]$, which lies well inside
the margin of $\pm 1$ percentage point. On this benchmark the pipeline
was thus marginally worse than the single agent on average, the reverse
of the exploratory result, but in both runs the difference is small
enough to count as equivalent. For H2b and H2c, the pipeline used
13\,310 input tokens per run against 19\,189 for the single agent, its API
cost was 21\,\% lower, and its latency was 20.4 seconds against 13.4
seconds.

The confirmatory estimates are close to the exploratory ones:
input tokens fell by 31\,\% in both runs, API cost by 21\,\% in both, and
latency rose by 52\,\% here against 69\,\% before. The secondary results
also replicate. The equivalence test is again not significant at the
narrower margin of $\pm 0.5$ percentage points in its $t$ version
($p = 0.11$), although its Wilcoxon version is. The smart timer again
outperforms both agent designs, with a mean gap of 0.17\,\% (Wilcoxon
$p < 0.001$ against each design, $r = 0.69$ for the single agent and
$0.64$ for the pipeline). Type D remains the hardest scenario type, with
mean gaps of 3.9\,\% for the single agent and 5.5\,\% for the pipeline,
while the LLM without tools fails most clearly in type E, with a mean gap
of 56\,\%.

The real home energy management tools were added after the
pre-registration, so the following comparisons are secondary and not
registered; their $p$-values are descriptive and not adjusted for
multiple testing. EMHASS as deployed reaches a mean gap of 0.36\,\% with
measured profiles, not significantly different from either agent design
(Wilcoxon $p = 0.90$ for the single agent and $p = 0.69$ for the
pipeline), and 3.5\,\% with nominal power, significantly worse than both
($p < 0.001$). The cheapest-hours automation exceeds the peak limit in 26
of the 48 scenarios, with a mean gap of 18.7\,\%. These patterns match
those on the main benchmark.

The registered findings, H1 for SQ1 and H2 for SQ2, are therefore not
specific to the benchmark on which the pipeline was developed: they hold
on unseen days with a different price resolution. They cover GPT-4o-mini
only. The registered hypothesis on a small model run locally (H3) was
withdrawn in the first amendment of the pre-registration, before its run
on the first benchmark was completed.

% ── References cited only in Appendix A: move back into the References list ──
% (they are inside \begin{hangparas} ... \end{hangparas} in main.tex)
%
% Department of Energy \& Climate Change. (2016). Domestic energy bills in
% 2015: The impact of variable consumption. In \textit{Energy Trends:
% March 2016}. UK Government. Retrieved 6 October 2026, from
% \url{https://assets.publishing.service.gov.uk/media/5a7f857740f0b6230268ffdc/Domestic_energy_bills_in_2015_-_The_impact_of_variable_consumption.pdf}
% 
% European Central Bank. (2026). \textit{Euro foreign exchange reference
% rates: Pound sterling (GBP)} [Data set]. Retrieved 6 October 2026, from
% \url{https://data.ecb.europa.eu/data/datasets/EXR}
  on the line before  \end{document}.
%   2. Move the two reference entries at the end of this file back into the References list of main.tex
%      (alphabetical order: after "Department of ..." entries / before "European Commission ..." etc.).
%   3. Restore the cross-references in the main text: search this file's labels app:uk2015 and app:earlier
%      (and sec:statistics, sec:baselines, sec:confirmatory, tab:baselines, tab:confirmatory,
%      tab:equivalence, tab:efficiency-tests, tab:uk-systems) and point the main text to them again
%      where it now says "the project repository" or names the results directly.
%
% Sources: results/uk2015.jsonl, results/tools/scenarios_uk2015.jsonl (Appendix A); results/runs.jsonl,
% results/baselines.jsonl, results/confirmation.jsonl, results/confirmation/confirmatory_tests.csv (Appendix B).

% ════════════════════════════════════════════════════════════════════════════
\appendix
\renewcommand{\thechapter}{\Alph{chapter}}

\chapter{Robustness check: UK prices of the REFIT period}
\label{app:uk2015}

% Data: data/scenarios_uk2015 (manifest), data/raw/prices_gb_entsoe.csv; results:
% results/uk2015.jsonl (configs/uk2015.json) and results/tools/scenarios_uk2015.jsonl (EMHASS).
The benchmarks of this thesis combine appliance cycles recorded in UK
homes in 2013--2015 with Finnish prices of 2022--2026
(Section~\ref{sec:dataset}). This appendix checks whether the results
also hold when the appliance data are priced in their own country and
period. It is a robustness check that was not pre-registered, and its
results are exploratory.

\section{Prices, tariff and scenarios}

The ENTSO-E Transparency Platform publishes day-ahead prices for the
Great Britain bidding zone only from 17 December 2014, in GBP per MWh
(ENTSO-E, 2026). This covers the last months of the REFIT recordings,
which ended in mid-2015 (Murray et al., 2017). The prices of 18 December
2014 to 31 July 2015, hourly, give 224 complete days on London time.
They were converted to euros at the mean of the European Central Bank's
daily reference rates for the same period, 0.73161 GBP per EUR
(European Central Bank, 2026), so that the agents' tools and the costs
use the same currency as in the other benchmarks. The saving in percent
does not depend on the currency.

The consumer price follows the form of the Finnish tariff
(Section~\ref{sec:dataset}): the spot price plus a constant charge per
kilowatt-hour, with VAT. The constant charge stands for all
non-wholesale costs of a UK household in 2015, that is network costs,
policy costs, the supplier's costs and margin, and the standing charge
spread over consumption. It was set so that the average consumer price
equals the average UK unit cost of standard electricity in 2015, 15.4
pence per kWh including 5\,\% VAT, which follows from DECC's average bill
of \pounds 584 for an annual consumption of 3\,800\,kWh (Department of
Energy \& Climate Change, 2016). With a mean spot price of 4.05 pence per
kWh in the window, the charge is 10.62 pence per kWh, or 0.1452 euros
per kWh. In 2014--2015 almost no UK household had a dynamic hourly
tariff, so the set describes a hypothetical household on such a tariff
rather than a real one.

The scenarios follow the factorial design of
Section~\ref{sec:factorial} with one change. The window contains no day
with a negative price, so the set has three price patterns, stable,
volatile and spike, each defined relative to the days of the window.
With two household loads, two peak limits and four days per pattern and
load, the set has 48 scenarios, 14 of which require coordination. The
cost of starting every appliance when it is ready has a median of
\EUR{0.63} per scenario. Every system of the confirmatory evaluation was
run on the set in the same way as on Q2: the Multi-Agent Pipeline and
the LLM without tools three times per scenario with GPT-4o-mini, and
EMHASS and the rule-based systems once.

\section{Results}

Table~\ref{tab:uk-systems} summarises the results. The ranking of the
systems is the same as on the Finnish benchmarks. The pipeline produced
a valid schedule in all 144 runs and saved 6.69\,\% of the cost without
scheduling, the same as the optimal schedule. EMHASS saved 6.45\,\% and
the smart timer 6.55\,\%. The cheapest-hours automations saved 2.6 to
3.6\,\%, and on the 14 scenarios that require coordination only 7\,\% of
their schedules were valid, against 100\,\% for the pipeline, EMHASS and
the smart timer. The LLM without tools saved 0.18\,\%, and a quarter of
its runs cost more than starting every appliance when ready. The
pipeline's mean latency was 18.5 seconds, and its API cost \$0.0025 per
recommendation.

\begin{table}[h!]
\caption{Robustness set: UK prices and tariff, December 2014 to July 2015 (48 scenarios; saving pooled over all runs).}
\label{tab:uk-systems}
\small
\begin{tabularx}{\textwidth}{>{\raggedright\arraybackslash}Xcccc}
\toprule
System & Valid & Saving & \makecell{Saving,\\strict limit} & \makecell{Saving,\\realistic limit} \\
\midrule
Multi-Agent Pipeline & 100\,\% & 6.69\,\% & 6.4\,\% & 6.9\,\% \\
Smart timer (own heuristic) & 100\,\% & 6.55\,\% & 6.2\,\% & 6.9\,\% \\
EMHASS, measured profiles & 100\,\% & 6.45\,\% & 6.0\,\% & 6.9\,\% \\
EMHASS, nominal power & 100\,\% & 6.06\,\% & 5.6\,\% & 6.5\,\% \\
Cheapest window & 70.8\,\% & 3.55\,\% & 1.5\,\% & 5.6\,\% \\
Cheapest-hours automation & 68.8\,\% & 2.62\,\% & 0.9\,\% & 4.3\,\% \\
LLM without tools & 95.1\,\% & 0.18\,\% & 0.3\,\% & 0.1\,\% \\
\midrule
Optimum (upper bound) & 100\,\% & 6.69\,\% & 6.4\,\% & 6.9\,\% \\
\bottomrule
\end{tabularx}
\end{table}

Paired per scenario, the pipeline saved more than the cheapest-hours
automation (by 3.3 percentage points), the cheapest window (2.1) and the
LLM without tools (6.4), with Wilcoxon $p < 0.001$ in each case, and it
saved the same as EMHASS within $\pm 1$ percentage point (mean difference
$+0.20$ points, 90\,\% confidence interval $[0.06, 0.34]$). These tests
mirror the registered ones but were not registered for this set, so they
are reported without adjustment.

The savings are about a quarter of those on the Finnish benchmarks.
Because the price patterns are defined relative to the days of the
window, even the volatile and spike days of 2015 are mild compared with
the Finnish ones: on these days the spot price varied within the day by
8 to 9 cents per kWh on average, against 16 to 20 cents on the
corresponding days of Q2. The window also has no negative prices. In addition,
the constant charge makes up most of the consumer price, so a given
shift in time changes a smaller share of the cost. The order of the
price patterns also differs from the Finnish sets: the stable days gave
the largest saving (9.8\,\% at the optimum) and the volatile days the
smallest (3.0\,\%). The reason was not analysed further.

\section{Conclusion}

Priced at the prices and tariff of their own country and period, the
REFIT appliance cycles give much smaller savings, but the comparison
between the systems does not change. The pipeline again reaches the
optimal saving with a valid schedule every time, matches EMHASS, and
saves more than the cheapest-hours automations, whose schedules again
fail where the household's appliances must be staggered. The main
results of this thesis therefore do not depend on combining UK
appliance data with Finnish prices, although the size of the saving
does depend on the price conditions.


% ════════════════════════════════════════════════════════════════════════════
\chapter{Evaluation of the earlier design}
\label{app:earlier}

This appendix reports the evaluation of the earlier design of the study,
summarised in Section~\ref{sec:earlier-design}. It compared the base
multi-agent pipeline, without retrieval or the optimiser tool, with a
single ReAct agent that had the first seven tools of
Table~\ref{tab:tools} and the same step budget, on the first benchmark of
48 scenarios in six types (Section~\ref{sec:dataset}) and on the
held-out Q set (Section~\ref{sec:resolution}). Its outcome was the
optimality gap, the cost of a schedule above the optimal cost in percent
of the optimal cost, with an infeasible schedule scored at the cost of
starting every appliance when ready. The research questions and
hypotheses of this design were replaced in the redesign, so the labels
SQ1, SQ2, H1 and H2 in Section~\ref{sec:confirmatory} refer to the first
pre-registration. The results remain valid for that design.

\section{Pipeline and single agent on the first benchmark}
\label{sec:statistics}

% Hand-typed from `hems analyze` (results/tables/equivalence.csv, efficiency_tests.csv), GPT-4o-mini only.
This section reports the tests of the earlier design for GPT-4o-mini on
the first benchmark, with the mean over three runs per scenario as the
paired observation ($N = 48$). These results are exploratory; the
pre-registered confirmatory results are reported in
Section~\ref{sec:confirmatory}. A run with the local model gpt-oss-20b
was stopped before completion and is not reported.

\paragraph{Schedule quality.} The multi-agent pipeline reached a mean
optimality gap of 1.25\,\% and the single agent 1.43\,\%. The Wilcoxon
signed-rank test finds no significant difference ($W = 449$,
$Z = -0.14$, $p = 0.89$, $r = 0.02$), and the median paired difference
is zero, because both designs reach the same cost in most scenarios. The
equivalence test shows that this is not merely a lack of evidence. The
mean paired difference is $-0.18$ percentage points, with a 90\,\%
confidence interval of $[-0.71, 0.35]$, which lies inside the primary
margin of $\pm 1$ percentage point (Table~\ref{tab:equivalence}). The two
designs are therefore equivalent in schedule quality within one
percentage point of the optimum. At the narrower margin of 0.5
percentage points the $t$ version of the test does not conclude
equivalence, because a few scenarios with large differences in either
direction widen the interval, although the Wilcoxon version does.

\begin{table}[h!]
\caption{Equivalence test (TOST) of the optimality gap, multi-agent vs.\ single-agent (GPT-4o-mini, $N = 48$).}
\label{tab:equivalence}
\small
\begin{tabular*}{\textwidth}{@{\extracolsep{\fill}}lccc}
\toprule
Margin & $p$, paired $t$ TOST & $p$, Wilcoxon TOST & Equivalent \\
\midrule
$\pm 0.5$ percentage points & 0.156 & $<$\,0.001 & No \\
$\pm 1$ percentage point (primary) & 0.006 & $<$\,0.001 & Yes \\
$\pm 2$ percentage points & $<$\,0.001 & $<$\,0.001 & Yes \\
\bottomrule
\end{tabular*}
\end{table}

\paragraph{Efficiency.} The two designs differ significantly in every
efficiency measure (Table~\ref{tab:efficiency-tests}). The pipeline used
34\,\% fewer input tokens, because each specialist works with a short
context and only its own tools, while the single agent carries the
entire conversation and all eight tool descriptions through every step.
It produced 83\,\% more output tokens and made 23\,\% more tool calls,
since four agents each write their own structured output and call their
tools separately. Input tokens dominate the price, so the API cost per run was
21\,\% lower for the pipeline, and lower in 35 of the 48 scenarios. The
sequential hand-offs made the pipeline 69\,\% slower, with a mean of
21.4 seconds against 12.7 seconds, and it was faster in only two
scenarios.

\begin{table}[h!]
\caption{Efficiency per run, multi-agent vs.\ single-agent (GPT-4o-mini; tokens and tool calls per run, API cost in USD, latency in seconds; each value is the mean over $N = 48$ scenarios of the per-scenario mean of three runs; Wilcoxon signed-rank test).}
\label{tab:efficiency-tests}
\small
\begin{tabular*}{\textwidth}{@{\extracolsep{\fill}}lcccccc}
\toprule
Measure & \makecell{Single\\agent} & \makecell{Multi-\\agent} & Difference & \makecell{Multi-agent\\lower} & $p$ & $r$ \\
\midrule
Input tokens & 21\,686 & 14\,292 & $-34$\,\% & 43/48 & $<$\,0.001 & 0.79 \\
Output tokens & 661 & 1\,208 & $+83$\,\% & 1/48 & $<$\,0.001 & 0.87 \\
Total tokens & 22\,347 & 15\,500 & $-31$\,\% & 40/48 & $<$\,0.001 & 0.75 \\
API cost (USD) & 0.0037 & 0.0029 & $-21$\,\% & 35/48 & $<$\,0.001 & 0.54 \\
Latency (s) & 12.7 & 21.4 & $+69$\,\% & 2/48 & $<$\,0.001 & 0.86 \\
Tool calls & 16.5 & 20.3 & $+23$\,\% & 1/48 & $<$\,0.001 & 0.81 \\
\bottomrule
\end{tabular*}
\end{table}

Taken together, decomposing the task into specialist agents did not
lower the cost of the schedules, but it reached the same quality with
about a third fewer tokens and a fifth lower inference cost, at the price
of a longer response time.

\section{Comparison with baselines and real tools}
\label{sec:baselines}

% Hand-typed from `hems analyze --runs results/runs.jsonl results/baselines.jsonl` (GPT-4o-mini).
Table~\ref{tab:baselines} compares the two agent designs with the smart
timer, the LLM without tools and the two real home energy management
tools. All systems that use a language model use GPT-4o-mini.

% EMHASS rows: Docker/REST service with default settings (tools/emhass_rest.py, systems emhass_rest*); cheapest hours;
% all in results/tools/scenarios.jsonl (secondary, added after the pre-registration). Savings captured = pooled.
\begin{table}[h!]
\caption{Agent designs, simpler baselines and real home energy management tools (GPT-4o-mini, main benchmark, 48 scenarios).}
\label{tab:baselines}
\small
\begin{tabular*}{\textwidth}{@{\extracolsep{\fill}}lcccccc}
\toprule
System & \makecell{Feasible\\runs} & \makecell{Mean\\gap} & \makecell{Median\\gap} & \makecell{Gap,\\type D} & \makecell{Savings\\captured} & \makecell{Time\\per run} \\
\midrule
Smart timer (own heuristic) & 48/48 & 0.23\,\% & 0.00\,\% & 1.2\,\% & 99\,\% & $<$\,0.01\,s \\
EMHASS, profiles & 48/48 & 0.32\,\% & 0.00\,\% & 0.1\,\% & 99\,\% & 1.2\,s \\
Multi-agent pipeline & 144/144 & 1.25\,\% & 0.01\,\% & 5.9\,\% & 93\,\% & 21.4\,s \\
Single agent & 143/144 & 1.43\,\% & 0.00\,\% & 6.5\,\% & 90\,\% & 12.7\,s \\
EMHASS, nominal power & 48/48 & 2.24\,\% & 0.45\,\% & 8.5\,\% & 87\,\% & 1.0\,s \\
Cheapest hours & 25/48 & 17.6\,\% & 4.2\,\% & 62.7\,\% & 25\,\% & $<$\,0.01\,s \\
LLM without tools & 126/144 & 23.1\,\% & 11.9\,\% & 66.0\,\% & 13\,\% & 2.0\,s \\
\bottomrule
\end{tabular*}
\end{table}

Without tools, GPT-4o-mini schedules poorly. Its mean optimality gap is
23.1\,\%, and 18 of its 144 schedules are infeasible: 14 exceed the peak
limit and 4 start a cycle outside its window. A further 22 of its
feasible schedules cost more than starting every appliance as soon as it
is ready, and over all its runs it captures only 13\,\% of the
achievable saving, against about 90\,\% for the agents. Both agent designs are significantly better (Wilcoxon
$p < 0.001$, $r = 0.87$ for each), with median paired differences of
10.8 percentage points for the single agent and 11.4 for the pipeline.
The same model that fails to schedule from a single prompt comes close to
the optimum when it can query prices, costs and loads through tools and
check its own work. Tool use, rather than the language model's own
calculation, is therefore what makes the agents viable.

The smart timer, however, outperforms both agent designs. It finds a
feasible schedule in every scenario, is optimal in 30 of the 48, and its
largest gap is 3.7\,\% (D04). Both agent designs have significantly
larger gaps than the timer (single agent $p < 0.001$, $r = 0.63$;
pipeline $p < 0.001$, $r = 0.68$). The difference is concentrated in
type D, where the timer's mean gap is 1.2\,\% against 5.9\,\% for the
pipeline and 6.5\,\% for the single agent. The median paired differences
are small (0.002 percentage points for the single agent and 0.017 for the
pipeline), because in most scenarios all three systems reach the optimum
or come very close to it. The timer needs no language model, runs in
about two milliseconds and has no inference cost. In this benchmark,
where the power profile and window of every appliance are known exactly
and the horizon is one day, placing the appliances one at a time is
almost always enough. The agents come within about one percentage point
of this rule but do not improve on it. What they add is the ability to
start from a request written in natural language rather than from
structured appliance data, and to explain the schedule they produce.

The two real home energy management tools, added after the
pre-registration, bracket the agents. Configured with the measured power
profiles, EMHASS comes close to the optimum: every schedule is feasible,
32 of the 48 are optimal, and the mean gap is 0.32\,\%. Its largest gap,
4.2\,\% in F02, comes from its default solver tolerance of 1\,\%
(the setting \texttt{lp\_solver\_mip\_rel\_gap} = 0.01; EMHASS
contributors, 2026), which
applies to the whole electricity bill including the base load and can
therefore amount to several percent of the appliance cost scored here.
With an exact tolerance the same optimization comes within 0.2\,\% of
the optimum in every scenario, which confirms that the MILP reference reflects what a
deployed open-source optimizer can achieve. The agents' mean gaps are
about one percentage point higher, but the difference from EMHASS as
deployed is not significant (Wilcoxon $p = 0.24$ for the single agent and
$p = 0.19$ for the pipeline). Configured with a constant nominal power,
EMHASS still keeps every schedule within the peak limit, because it
plans each cycle at its peak power, but it misjudges the cost of the real
cycles and reaches a mean gap of 2.2\,\%. Both agent designs are
significantly better than this configuration (Wilcoxon $p < 0.001$ for
each, $r = 0.67$ for the single agent and $0.64$ for the pipeline). The
cheapest-hours automation, under the assumption stated in
Section~\ref{sec:evaluation} that its trigger starts a cycle which then
runs to completion, exceeds the peak limit in 23 of the 48 scenarios,
because each appliance follows
its own cheap period and appliances whose windows share a price minimum
start together. Scored at the naive cost when infeasible, it reaches a
mean gap of 17.6\,\%, and both agent designs are significantly better
($p < 0.001$, $r = 0.82$ and $0.81$). The agents are therefore
outperformed by the smart timer and are not significantly different from
a well-configured optimizer as deployed, while they outperform a simply
configured optimizer and the per-device automation that is common among
households.

Two checks show that this ranking does not depend on how the costs are
scored. First, scoring an infeasible run at the cost of the naive
schedule does not drive the result. Restricted to feasible runs, the
mean gaps are 0.23\,\% for the smart timer, 0.32\,\% for EMHASS with
profiles, 1.25\,\% for the pipeline, 1.45\,\% for the single agent,
2.24\,\% for EMHASS with nominal power, 7.9\,\% for the cheapest-hours
automation and 23.0\,\% for the LLM without tools, the same order as in
Table~\ref{tab:baselines}, and the order on the held-out benchmark is
also unchanged. The cheapest-hours automation looks better on feasible
runs only because these are the days on which its appliances happen not
to collide. Second, the ranking is the same under the spot price alone
as under the consumer price, as explained in
Section~\ref{sec:evaluation}.

\section{Confirmatory run on the held-out Q set}
\label{sec:confirmatory}

% Hand-typed from `hems confirm` (results/confirmation/confirmatory_tests.csv) and
% `hems analyze --runs results/confirmation.jsonl` for the secondary results.
The hypotheses of the earlier SQ1 and SQ2 were fixed in the pre-registration, which
was committed to the project repository on 30 September 2026 together
with the code, prompts, tools and analysis script. The confirmatory run
followed on the held-out benchmark of 48 scenarios with fifteen-minute
prices (Section~\ref{sec:resolution}), on which no language model had
been run before. All systems that use a language model used GPT-4o-mini
with three runs per scenario, and the analysis was run exactly as
registered. The run cost USD~0.92 in API fees. Table~\ref{tab:confirmatory}
reports the six registered tests.

\begin{table}[h!]
\caption{Pre-registered confirmatory tests on the held-out fifteen-minute benchmark (GPT-4o-mini, $N = 48$; $p$-values Holm-adjusted within each sub-question).}
\label{tab:confirmatory}
\small
\begin{tabularx}{\textwidth}{l>{\raggedright\arraybackslash}Xcccc}
\toprule
& Hypothesis & Estimate & $p$ & $r$ & Supported \\
\midrule
H1 & Single agent gap $<$ LLM without tools & $-17.9$ pp & $<$\,0.001 & 0.86 & Yes \\
H1 & Pipeline gap $<$ LLM without tools & $-17.7$ pp & $<$\,0.001 & 0.86 & Yes \\
H2a & Pipeline equivalent to single agent ($\pm 1$ pp) & $+0.22$ pp & $<$\,0.001 & -- & Yes \\
H2b & Pipeline uses fewer input tokens & $-31$\,\% & $<$\,0.001 & 0.84 & Yes \\
H2b & Pipeline has a lower API cost per run & $-21$\,\% & $<$\,0.001 & 0.70 & Yes \\
H2c & Pipeline has a longer latency & $+52$\,\% & $<$\,0.001 & 0.80 & Yes \\
\bottomrule
\end{tabularx}
\end{table}

All six hypotheses are supported. For H1, both agent designs reached a
mean optimality gap of about one percent (0.86\,\% for the single agent
and 1.08\,\% for the pipeline), whereas the LLM without tools reached
18.8\,\%. Fifteen of its 144 schedules were infeasible, and over all its
runs it captured only 24\,\% of the achievable saving, against 92\,\% for
the single agent and 90\,\% for the pipeline. For H2a,
the mean paired difference in gap was $+0.22$ percentage points with a
90\,\% confidence interval of $[-0.15, 0.59]$, which lies well inside
the margin of $\pm 1$ percentage point. On this benchmark the pipeline
was thus marginally worse than the single agent on average, the reverse
of the exploratory result, but in both runs the difference is small
enough to count as equivalent. For H2b and H2c, the pipeline used
13\,310 input tokens per run against 19\,189 for the single agent, its API
cost was 21\,\% lower, and its latency was 20.4 seconds against 13.4
seconds.

The confirmatory estimates are close to the exploratory ones:
input tokens fell by 31\,\% in both runs, API cost by 21\,\% in both, and
latency rose by 52\,\% here against 69\,\% before. The secondary results
also replicate. The equivalence test is again not significant at the
narrower margin of $\pm 0.5$ percentage points in its $t$ version
($p = 0.11$), although its Wilcoxon version is. The smart timer again
outperforms both agent designs, with a mean gap of 0.17\,\% (Wilcoxon
$p < 0.001$ against each design, $r = 0.69$ for the single agent and
$0.64$ for the pipeline). Type D remains the hardest scenario type, with
mean gaps of 3.9\,\% for the single agent and 5.5\,\% for the pipeline,
while the LLM without tools fails most clearly in type E, with a mean gap
of 56\,\%.

The real home energy management tools were added after the
pre-registration, so the following comparisons are secondary and not
registered; their $p$-values are descriptive and not adjusted for
multiple testing. EMHASS as deployed reaches a mean gap of 0.36\,\% with
measured profiles, not significantly different from either agent design
(Wilcoxon $p = 0.90$ for the single agent and $p = 0.69$ for the
pipeline), and 3.5\,\% with nominal power, significantly worse than both
($p < 0.001$). The cheapest-hours automation exceeds the peak limit in 26
of the 48 scenarios, with a mean gap of 18.7\,\%. These patterns match
those on the main benchmark.

The registered findings, H1 for SQ1 and H2 for SQ2, are therefore not
specific to the benchmark on which the pipeline was developed: they hold
on unseen days with a different price resolution. They cover GPT-4o-mini
only. The registered hypothesis on a small model run locally (H3) was
withdrawn in the first amendment of the pre-registration, before its run
on the first benchmark was completed.

% ── References cited only in Appendix A: move back into the References list ──
% (they are inside \begin{hangparas} ... \end{hangparas} in main.tex)
%
% Department of Energy \& Climate Change. (2016). Domestic energy bills in
% 2015: The impact of variable consumption. In \textit{Energy Trends:
% March 2016}. UK Government. Retrieved 6 October 2026, from
% \url{https://assets.publishing.service.gov.uk/media/5a7f857740f0b6230268ffdc/Domestic_energy_bills_in_2015_-_The_impact_of_variable_consumption.pdf}
% 
% European Central Bank. (2026). \textit{Euro foreign exchange reference
% rates: Pound sterling (GBP)} [Data set]. Retrieved 6 October 2026, from
% \url{https://data.ecb.europa.eu/data/datasets/EXR}
  on the line before  \end{document}.
%   2. Move the two reference entries at the end of this file back into the References list of main.tex
%      (alphabetical order: after "Department of ..." entries / before "European Commission ..." etc.).
%   3. Restore the cross-references in the main text: search this file's labels app:uk2015 and app:earlier
%      (and sec:statistics, sec:baselines, sec:confirmatory, tab:baselines, tab:confirmatory,
%      tab:equivalence, tab:efficiency-tests, tab:uk-systems) and point the main text to them again
%      where it now says "the project repository" or names the results directly.
%
% Sources: results/uk2015.jsonl, results/tools/scenarios_uk2015.jsonl (Appendix A); results/runs.jsonl,
% results/baselines.jsonl, results/confirmation.jsonl, results/confirmation/confirmatory_tests.csv (Appendix B).

% ════════════════════════════════════════════════════════════════════════════
\appendix
\renewcommand{\thechapter}{\Alph{chapter}}

\chapter{Robustness check: UK prices of the REFIT period}
\label{app:uk2015}

% Data: data/scenarios_uk2015 (manifest), data/raw/prices_gb_entsoe.csv; results:
% results/uk2015.jsonl (configs/uk2015.json) and results/tools/scenarios_uk2015.jsonl (EMHASS).
The benchmarks of this thesis combine appliance cycles recorded in UK
homes in 2013--2015 with Finnish prices of 2022--2026
(Section~\ref{sec:dataset}). This appendix checks whether the results
also hold when the appliance data are priced in their own country and
period. It is a robustness check that was not pre-registered, and its
results are exploratory.

\section{Prices, tariff and scenarios}

The ENTSO-E Transparency Platform publishes day-ahead prices for the
Great Britain bidding zone only from 17 December 2014, in GBP per MWh
(ENTSO-E, 2026). This covers the last months of the REFIT recordings,
which ended in mid-2015 (Murray et al., 2017). The prices of 18 December
2014 to 31 July 2015, hourly, give 224 complete days on London time.
They were converted to euros at the mean of the European Central Bank's
daily reference rates for the same period, 0.73161 GBP per EUR
(European Central Bank, 2026), so that the agents' tools and the costs
use the same currency as in the other benchmarks. The saving in percent
does not depend on the currency.

The consumer price follows the form of the Finnish tariff
(Section~\ref{sec:dataset}): the spot price plus a constant charge per
kilowatt-hour, with VAT. The constant charge stands for all
non-wholesale costs of a UK household in 2015, that is network costs,
policy costs, the supplier's costs and margin, and the standing charge
spread over consumption. It was set so that the average consumer price
equals the average UK unit cost of standard electricity in 2015, 15.4
pence per kWh including 5\,\% VAT, which follows from DECC's average bill
of \pounds 584 for an annual consumption of 3\,800\,kWh (Department of
Energy \& Climate Change, 2016). With a mean spot price of 4.05 pence per
kWh in the window, the charge is 10.62 pence per kWh, or 0.1452 euros
per kWh. In 2014--2015 almost no UK household had a dynamic hourly
tariff, so the set describes a hypothetical household on such a tariff
rather than a real one.

The scenarios follow the factorial design of
Section~\ref{sec:factorial} with one change. The window contains no day
with a negative price, so the set has three price patterns, stable,
volatile and spike, each defined relative to the days of the window.
With two household loads, two peak limits and four days per pattern and
load, the set has 48 scenarios, 14 of which require coordination. The
cost of starting every appliance when it is ready has a median of
\EUR{0.63} per scenario. Every system of the confirmatory evaluation was
run on the set in the same way as on Q2: the Multi-Agent Pipeline and
the LLM without tools three times per scenario with GPT-4o-mini, and
EMHASS and the rule-based systems once.

\section{Results}

Table~\ref{tab:uk-systems} summarises the results. The ranking of the
systems is the same as on the Finnish benchmarks. The pipeline produced
a valid schedule in all 144 runs and saved 6.69\,\% of the cost without
scheduling, the same as the optimal schedule. EMHASS saved 6.45\,\% and
the smart timer 6.55\,\%. The cheapest-hours automations saved 2.6 to
3.6\,\%, and on the 14 scenarios that require coordination only 7\,\% of
their schedules were valid, against 100\,\% for the pipeline, EMHASS and
the smart timer. The LLM without tools saved 0.18\,\%, and a quarter of
its runs cost more than starting every appliance when ready. The
pipeline's mean latency was 18.5 seconds, and its API cost \$0.0025 per
recommendation.

\begin{table}[h!]
\caption{Robustness set: UK prices and tariff, December 2014 to July 2015 (48 scenarios; saving pooled over all runs).}
\label{tab:uk-systems}
\small
\begin{tabularx}{\textwidth}{>{\raggedright\arraybackslash}Xcccc}
\toprule
System & Valid & Saving & \makecell{Saving,\\strict limit} & \makecell{Saving,\\realistic limit} \\
\midrule
Multi-Agent Pipeline & 100\,\% & 6.69\,\% & 6.4\,\% & 6.9\,\% \\
Smart timer (own heuristic) & 100\,\% & 6.55\,\% & 6.2\,\% & 6.9\,\% \\
EMHASS, measured profiles & 100\,\% & 6.45\,\% & 6.0\,\% & 6.9\,\% \\
EMHASS, nominal power & 100\,\% & 6.06\,\% & 5.6\,\% & 6.5\,\% \\
Cheapest window & 70.8\,\% & 3.55\,\% & 1.5\,\% & 5.6\,\% \\
Cheapest-hours automation & 68.8\,\% & 2.62\,\% & 0.9\,\% & 4.3\,\% \\
LLM without tools & 95.1\,\% & 0.18\,\% & 0.3\,\% & 0.1\,\% \\
\midrule
Optimum (upper bound) & 100\,\% & 6.69\,\% & 6.4\,\% & 6.9\,\% \\
\bottomrule
\end{tabularx}
\end{table}

Paired per scenario, the pipeline saved more than the cheapest-hours
automation (by 3.3 percentage points), the cheapest window (2.1) and the
LLM without tools (6.4), with Wilcoxon $p < 0.001$ in each case, and it
saved the same as EMHASS within $\pm 1$ percentage point (mean difference
$+0.20$ points, 90\,\% confidence interval $[0.06, 0.34]$). These tests
mirror the registered ones but were not registered for this set, so they
are reported without adjustment.

The savings are about a quarter of those on the Finnish benchmarks.
Because the price patterns are defined relative to the days of the
window, even the volatile and spike days of 2015 are mild compared with
the Finnish ones: on these days the spot price varied within the day by
8 to 9 cents per kWh on average, against 16 to 20 cents on the
corresponding days of Q2. The window also has no negative prices. In addition,
the constant charge makes up most of the consumer price, so a given
shift in time changes a smaller share of the cost. The order of the
price patterns also differs from the Finnish sets: the stable days gave
the largest saving (9.8\,\% at the optimum) and the volatile days the
smallest (3.0\,\%). The reason was not analysed further.

\section{Conclusion}

Priced at the prices and tariff of their own country and period, the
REFIT appliance cycles give much smaller savings, but the comparison
between the systems does not change. The pipeline again reaches the
optimal saving with a valid schedule every time, matches EMHASS, and
saves more than the cheapest-hours automations, whose schedules again
fail where the household's appliances must be staggered. The main
results of this thesis therefore do not depend on combining UK
appliance data with Finnish prices, although the size of the saving
does depend on the price conditions.


% ════════════════════════════════════════════════════════════════════════════
\chapter{Evaluation of the earlier design}
\label{app:earlier}

This appendix reports the evaluation of the earlier design of the study,
summarised in Section~\ref{sec:earlier-design}. It compared the base
multi-agent pipeline, without retrieval or the optimiser tool, with a
single ReAct agent that had the first seven tools of
Table~\ref{tab:tools} and the same step budget, on the first benchmark of
48 scenarios in six types (Section~\ref{sec:dataset}) and on the
held-out Q set (Section~\ref{sec:resolution}). Its outcome was the
optimality gap, the cost of a schedule above the optimal cost in percent
of the optimal cost, with an infeasible schedule scored at the cost of
starting every appliance when ready. The research questions and
hypotheses of this design were replaced in the redesign, so the labels
SQ1, SQ2, H1 and H2 in Section~\ref{sec:confirmatory} refer to the first
pre-registration. The results remain valid for that design.

\section{Pipeline and single agent on the first benchmark}
\label{sec:statistics}

% Hand-typed from `hems analyze` (results/tables/equivalence.csv, efficiency_tests.csv), GPT-4o-mini only.
This section reports the tests of the earlier design for GPT-4o-mini on
the first benchmark, with the mean over three runs per scenario as the
paired observation ($N = 48$). These results are exploratory; the
pre-registered confirmatory results are reported in
Section~\ref{sec:confirmatory}. A run with the local model gpt-oss-20b
was stopped before completion and is not reported.

\paragraph{Schedule quality.} The multi-agent pipeline reached a mean
optimality gap of 1.25\,\% and the single agent 1.43\,\%. The Wilcoxon
signed-rank test finds no significant difference ($W = 449$,
$Z = -0.14$, $p = 0.89$, $r = 0.02$), and the median paired difference
is zero, because both designs reach the same cost in most scenarios. The
equivalence test shows that this is not merely a lack of evidence. The
mean paired difference is $-0.18$ percentage points, with a 90\,\%
confidence interval of $[-0.71, 0.35]$, which lies inside the primary
margin of $\pm 1$ percentage point (Table~\ref{tab:equivalence}). The two
designs are therefore equivalent in schedule quality within one
percentage point of the optimum. At the narrower margin of 0.5
percentage points the $t$ version of the test does not conclude
equivalence, because a few scenarios with large differences in either
direction widen the interval, although the Wilcoxon version does.

\begin{table}[h!]
\caption{Equivalence test (TOST) of the optimality gap, multi-agent vs.\ single-agent (GPT-4o-mini, $N = 48$).}
\label{tab:equivalence}
\small
\begin{tabular*}{\textwidth}{@{\extracolsep{\fill}}lccc}
\toprule
Margin & $p$, paired $t$ TOST & $p$, Wilcoxon TOST & Equivalent \\
\midrule
$\pm 0.5$ percentage points & 0.156 & $<$\,0.001 & No \\
$\pm 1$ percentage point (primary) & 0.006 & $<$\,0.001 & Yes \\
$\pm 2$ percentage points & $<$\,0.001 & $<$\,0.001 & Yes \\
\bottomrule
\end{tabular*}
\end{table}

\paragraph{Efficiency.} The two designs differ significantly in every
efficiency measure (Table~\ref{tab:efficiency-tests}). The pipeline used
34\,\% fewer input tokens, because each specialist works with a short
context and only its own tools, while the single agent carries the
entire conversation and all eight tool descriptions through every step.
It produced 83\,\% more output tokens and made 23\,\% more tool calls,
since four agents each write their own structured output and call their
tools separately. Input tokens dominate the price, so the API cost per run was
21\,\% lower for the pipeline, and lower in 35 of the 48 scenarios. The
sequential hand-offs made the pipeline 69\,\% slower, with a mean of
21.4 seconds against 12.7 seconds, and it was faster in only two
scenarios.

\begin{table}[h!]
\caption{Efficiency per run, multi-agent vs.\ single-agent (GPT-4o-mini; tokens and tool calls per run, API cost in USD, latency in seconds; each value is the mean over $N = 48$ scenarios of the per-scenario mean of three runs; Wilcoxon signed-rank test).}
\label{tab:efficiency-tests}
\small
\begin{tabular*}{\textwidth}{@{\extracolsep{\fill}}lcccccc}
\toprule
Measure & \makecell{Single\\agent} & \makecell{Multi-\\agent} & Difference & \makecell{Multi-agent\\lower} & $p$ & $r$ \\
\midrule
Input tokens & 21\,686 & 14\,292 & $-34$\,\% & 43/48 & $<$\,0.001 & 0.79 \\
Output tokens & 661 & 1\,208 & $+83$\,\% & 1/48 & $<$\,0.001 & 0.87 \\
Total tokens & 22\,347 & 15\,500 & $-31$\,\% & 40/48 & $<$\,0.001 & 0.75 \\
API cost (USD) & 0.0037 & 0.0029 & $-21$\,\% & 35/48 & $<$\,0.001 & 0.54 \\
Latency (s) & 12.7 & 21.4 & $+69$\,\% & 2/48 & $<$\,0.001 & 0.86 \\
Tool calls & 16.5 & 20.3 & $+23$\,\% & 1/48 & $<$\,0.001 & 0.81 \\
\bottomrule
\end{tabular*}
\end{table}

Taken together, decomposing the task into specialist agents did not
lower the cost of the schedules, but it reached the same quality with
about a third fewer tokens and a fifth lower inference cost, at the price
of a longer response time.

\section{Comparison with baselines and real tools}
\label{sec:baselines}

% Hand-typed from `hems analyze --runs results/runs.jsonl results/baselines.jsonl` (GPT-4o-mini).
Table~\ref{tab:baselines} compares the two agent designs with the smart
timer, the LLM without tools and the two real home energy management
tools. All systems that use a language model use GPT-4o-mini.

% EMHASS rows: Docker/REST service with default settings (tools/emhass_rest.py, systems emhass_rest*); cheapest hours;
% all in results/tools/scenarios.jsonl (secondary, added after the pre-registration). Savings captured = pooled.
\begin{table}[h!]
\caption{Agent designs, simpler baselines and real home energy management tools (GPT-4o-mini, main benchmark, 48 scenarios).}
\label{tab:baselines}
\small
\begin{tabular*}{\textwidth}{@{\extracolsep{\fill}}lcccccc}
\toprule
System & \makecell{Feasible\\runs} & \makecell{Mean\\gap} & \makecell{Median\\gap} & \makecell{Gap,\\type D} & \makecell{Savings\\captured} & \makecell{Time\\per run} \\
\midrule
Smart timer (own heuristic) & 48/48 & 0.23\,\% & 0.00\,\% & 1.2\,\% & 99\,\% & $<$\,0.01\,s \\
EMHASS, profiles & 48/48 & 0.32\,\% & 0.00\,\% & 0.1\,\% & 99\,\% & 1.2\,s \\
Multi-agent pipeline & 144/144 & 1.25\,\% & 0.01\,\% & 5.9\,\% & 93\,\% & 21.4\,s \\
Single agent & 143/144 & 1.43\,\% & 0.00\,\% & 6.5\,\% & 90\,\% & 12.7\,s \\
EMHASS, nominal power & 48/48 & 2.24\,\% & 0.45\,\% & 8.5\,\% & 87\,\% & 1.0\,s \\
Cheapest hours & 25/48 & 17.6\,\% & 4.2\,\% & 62.7\,\% & 25\,\% & $<$\,0.01\,s \\
LLM without tools & 126/144 & 23.1\,\% & 11.9\,\% & 66.0\,\% & 13\,\% & 2.0\,s \\
\bottomrule
\end{tabular*}
\end{table}

Without tools, GPT-4o-mini schedules poorly. Its mean optimality gap is
23.1\,\%, and 18 of its 144 schedules are infeasible: 14 exceed the peak
limit and 4 start a cycle outside its window. A further 22 of its
feasible schedules cost more than starting every appliance as soon as it
is ready, and over all its runs it captures only 13\,\% of the
achievable saving, against about 90\,\% for the agents. Both agent designs are significantly better (Wilcoxon
$p < 0.001$, $r = 0.87$ for each), with median paired differences of
10.8 percentage points for the single agent and 11.4 for the pipeline.
The same model that fails to schedule from a single prompt comes close to
the optimum when it can query prices, costs and loads through tools and
check its own work. Tool use, rather than the language model's own
calculation, is therefore what makes the agents viable.

The smart timer, however, outperforms both agent designs. It finds a
feasible schedule in every scenario, is optimal in 30 of the 48, and its
largest gap is 3.7\,\% (D04). Both agent designs have significantly
larger gaps than the timer (single agent $p < 0.001$, $r = 0.63$;
pipeline $p < 0.001$, $r = 0.68$). The difference is concentrated in
type D, where the timer's mean gap is 1.2\,\% against 5.9\,\% for the
pipeline and 6.5\,\% for the single agent. The median paired differences
are small (0.002 percentage points for the single agent and 0.017 for the
pipeline), because in most scenarios all three systems reach the optimum
or come very close to it. The timer needs no language model, runs in
about two milliseconds and has no inference cost. In this benchmark,
where the power profile and window of every appliance are known exactly
and the horizon is one day, placing the appliances one at a time is
almost always enough. The agents come within about one percentage point
of this rule but do not improve on it. What they add is the ability to
start from a request written in natural language rather than from
structured appliance data, and to explain the schedule they produce.

The two real home energy management tools, added after the
pre-registration, bracket the agents. Configured with the measured power
profiles, EMHASS comes close to the optimum: every schedule is feasible,
32 of the 48 are optimal, and the mean gap is 0.32\,\%. Its largest gap,
4.2\,\% in F02, comes from its default solver tolerance of 1\,\%
(the setting \texttt{lp\_solver\_mip\_rel\_gap} = 0.01; EMHASS
contributors, 2026), which
applies to the whole electricity bill including the base load and can
therefore amount to several percent of the appliance cost scored here.
With an exact tolerance the same optimization comes within 0.2\,\% of
the optimum in every scenario, which confirms that the MILP reference reflects what a
deployed open-source optimizer can achieve. The agents' mean gaps are
about one percentage point higher, but the difference from EMHASS as
deployed is not significant (Wilcoxon $p = 0.24$ for the single agent and
$p = 0.19$ for the pipeline). Configured with a constant nominal power,
EMHASS still keeps every schedule within the peak limit, because it
plans each cycle at its peak power, but it misjudges the cost of the real
cycles and reaches a mean gap of 2.2\,\%. Both agent designs are
significantly better than this configuration (Wilcoxon $p < 0.001$ for
each, $r = 0.67$ for the single agent and $0.64$ for the pipeline). The
cheapest-hours automation, under the assumption stated in
Section~\ref{sec:evaluation} that its trigger starts a cycle which then
runs to completion, exceeds the peak limit in 23 of the 48 scenarios,
because each appliance follows
its own cheap period and appliances whose windows share a price minimum
start together. Scored at the naive cost when infeasible, it reaches a
mean gap of 17.6\,\%, and both agent designs are significantly better
($p < 0.001$, $r = 0.82$ and $0.81$). The agents are therefore
outperformed by the smart timer and are not significantly different from
a well-configured optimizer as deployed, while they outperform a simply
configured optimizer and the per-device automation that is common among
households.

Two checks show that this ranking does not depend on how the costs are
scored. First, scoring an infeasible run at the cost of the naive
schedule does not drive the result. Restricted to feasible runs, the
mean gaps are 0.23\,\% for the smart timer, 0.32\,\% for EMHASS with
profiles, 1.25\,\% for the pipeline, 1.45\,\% for the single agent,
2.24\,\% for EMHASS with nominal power, 7.9\,\% for the cheapest-hours
automation and 23.0\,\% for the LLM without tools, the same order as in
Table~\ref{tab:baselines}, and the order on the held-out benchmark is
also unchanged. The cheapest-hours automation looks better on feasible
runs only because these are the days on which its appliances happen not
to collide. Second, the ranking is the same under the spot price alone
as under the consumer price, as explained in
Section~\ref{sec:evaluation}.

\section{Confirmatory run on the held-out Q set}
\label{sec:confirmatory}

% Hand-typed from `hems confirm` (results/confirmation/confirmatory_tests.csv) and
% `hems analyze --runs results/confirmation.jsonl` for the secondary results.
The hypotheses of the earlier SQ1 and SQ2 were fixed in the pre-registration, which
was committed to the project repository on 30 September 2026 together
with the code, prompts, tools and analysis script. The confirmatory run
followed on the held-out benchmark of 48 scenarios with fifteen-minute
prices (Section~\ref{sec:resolution}), on which no language model had
been run before. All systems that use a language model used GPT-4o-mini
with three runs per scenario, and the analysis was run exactly as
registered. The run cost USD~0.92 in API fees. Table~\ref{tab:confirmatory}
reports the six registered tests.

\begin{table}[h!]
\caption{Pre-registered confirmatory tests on the held-out fifteen-minute benchmark (GPT-4o-mini, $N = 48$; $p$-values Holm-adjusted within each sub-question).}
\label{tab:confirmatory}
\small
\begin{tabularx}{\textwidth}{l>{\raggedright\arraybackslash}Xcccc}
\toprule
& Hypothesis & Estimate & $p$ & $r$ & Supported \\
\midrule
H1 & Single agent gap $<$ LLM without tools & $-17.9$ pp & $<$\,0.001 & 0.86 & Yes \\
H1 & Pipeline gap $<$ LLM without tools & $-17.7$ pp & $<$\,0.001 & 0.86 & Yes \\
H2a & Pipeline equivalent to single agent ($\pm 1$ pp) & $+0.22$ pp & $<$\,0.001 & -- & Yes \\
H2b & Pipeline uses fewer input tokens & $-31$\,\% & $<$\,0.001 & 0.84 & Yes \\
H2b & Pipeline has a lower API cost per run & $-21$\,\% & $<$\,0.001 & 0.70 & Yes \\
H2c & Pipeline has a longer latency & $+52$\,\% & $<$\,0.001 & 0.80 & Yes \\
\bottomrule
\end{tabularx}
\end{table}

All six hypotheses are supported. For H1, both agent designs reached a
mean optimality gap of about one percent (0.86\,\% for the single agent
and 1.08\,\% for the pipeline), whereas the LLM without tools reached
18.8\,\%. Fifteen of its 144 schedules were infeasible, and over all its
runs it captured only 24\,\% of the achievable saving, against 92\,\% for
the single agent and 90\,\% for the pipeline. For H2a,
the mean paired difference in gap was $+0.22$ percentage points with a
90\,\% confidence interval of $[-0.15, 0.59]$, which lies well inside
the margin of $\pm 1$ percentage point. On this benchmark the pipeline
was thus marginally worse than the single agent on average, the reverse
of the exploratory result, but in both runs the difference is small
enough to count as equivalent. For H2b and H2c, the pipeline used
13\,310 input tokens per run against 19\,189 for the single agent, its API
cost was 21\,\% lower, and its latency was 20.4 seconds against 13.4
seconds.

The confirmatory estimates are close to the exploratory ones:
input tokens fell by 31\,\% in both runs, API cost by 21\,\% in both, and
latency rose by 52\,\% here against 69\,\% before. The secondary results
also replicate. The equivalence test is again not significant at the
narrower margin of $\pm 0.5$ percentage points in its $t$ version
($p = 0.11$), although its Wilcoxon version is. The smart timer again
outperforms both agent designs, with a mean gap of 0.17\,\% (Wilcoxon
$p < 0.001$ against each design, $r = 0.69$ for the single agent and
$0.64$ for the pipeline). Type D remains the hardest scenario type, with
mean gaps of 3.9\,\% for the single agent and 5.5\,\% for the pipeline,
while the LLM without tools fails most clearly in type E, with a mean gap
of 56\,\%.

The real home energy management tools were added after the
pre-registration, so the following comparisons are secondary and not
registered; their $p$-values are descriptive and not adjusted for
multiple testing. EMHASS as deployed reaches a mean gap of 0.36\,\% with
measured profiles, not significantly different from either agent design
(Wilcoxon $p = 0.90$ for the single agent and $p = 0.69$ for the
pipeline), and 3.5\,\% with nominal power, significantly worse than both
($p < 0.001$). The cheapest-hours automation exceeds the peak limit in 26
of the 48 scenarios, with a mean gap of 18.7\,\%. These patterns match
those on the main benchmark.

The registered findings, H1 for SQ1 and H2 for SQ2, are therefore not
specific to the benchmark on which the pipeline was developed: they hold
on unseen days with a different price resolution. They cover GPT-4o-mini
only. The registered hypothesis on a small model run locally (H3) was
withdrawn in the first amendment of the pre-registration, before its run
on the first benchmark was completed.

% ── References cited only in Appendix A: move back into the References list ──
% (they are inside \begin{hangparas} ... \end{hangparas} in main.tex)
%
% Department of Energy \& Climate Change. (2016). Domestic energy bills in
% 2015: The impact of variable consumption. In \textit{Energy Trends:
% March 2016}. UK Government. Retrieved 6 October 2026, from
% \url{https://assets.publishing.service.gov.uk/media/5a7f857740f0b6230268ffdc/Domestic_energy_bills_in_2015_-_The_impact_of_variable_consumption.pdf}
% 
% European Central Bank. (2026). \textit{Euro foreign exchange reference
% rates: Pound sterling (GBP)} [Data set]. Retrieved 6 October 2026, from
% \url{https://data.ecb.europa.eu/data/datasets/EXR}
  on the line before  \end{document}.
%   2. Move the two reference entries at the end of this file back into the References list of main.tex
%      (alphabetical order: after "Department of ..." entries / before "European Commission ..." etc.).
%   3. Restore the cross-references in the main text: search this file's labels app:uk2015 and app:earlier
%      (and sec:statistics, sec:baselines, sec:confirmatory, tab:baselines, tab:confirmatory,
%      tab:equivalence, tab:efficiency-tests, tab:uk-systems) and point the main text to them again
%      where it now says "the project repository" or names the results directly.
%
% Sources: results/uk2015.jsonl, results/tools/scenarios_uk2015.jsonl (Appendix A); results/runs.jsonl,
% results/baselines.jsonl, results/confirmation.jsonl, results/confirmation/confirmatory_tests.csv (Appendix B).

% ════════════════════════════════════════════════════════════════════════════
\appendix
\renewcommand{\thechapter}{\Alph{chapter}}

\chapter{Robustness check: UK prices of the REFIT period}
\label{app:uk2015}

% Data: data/scenarios_uk2015 (manifest), data/raw/prices_gb_entsoe.csv; results:
% results/uk2015.jsonl (configs/uk2015.json) and results/tools/scenarios_uk2015.jsonl (EMHASS).
The benchmarks of this thesis combine appliance cycles recorded in UK
homes in 2013--2015 with Finnish prices of 2022--2026
(Section~\ref{sec:dataset}). This appendix checks whether the results
also hold when the appliance data are priced in their own country and
period. It is a robustness check that was not pre-registered, and its
results are exploratory.

\section{Prices, tariff and scenarios}

The ENTSO-E Transparency Platform publishes day-ahead prices for the
Great Britain bidding zone only from 17 December 2014, in GBP per MWh
(ENTSO-E, 2026). This covers the last months of the REFIT recordings,
which ended in mid-2015 (Murray et al., 2017). The prices of 18 December
2014 to 31 July 2015, hourly, give 224 complete days on London time.
They were converted to euros at the mean of the European Central Bank's
daily reference rates for the same period, 0.73161 GBP per EUR
(European Central Bank, 2026), so that the agents' tools and the costs
use the same currency as in the other benchmarks. The saving in percent
does not depend on the currency.

The consumer price follows the form of the Finnish tariff
(Section~\ref{sec:dataset}): the spot price plus a constant charge per
kilowatt-hour, with VAT. The constant charge stands for all
non-wholesale costs of a UK household in 2015, that is network costs,
policy costs, the supplier's costs and margin, and the standing charge
spread over consumption. It was set so that the average consumer price
equals the average UK unit cost of standard electricity in 2015, 15.4
pence per kWh including 5\,\% VAT, which follows from DECC's average bill
of \pounds 584 for an annual consumption of 3\,800\,kWh (Department of
Energy \& Climate Change, 2016). With a mean spot price of 4.05 pence per
kWh in the window, the charge is 10.62 pence per kWh, or 0.1452 euros
per kWh. In 2014--2015 almost no UK household had a dynamic hourly
tariff, so the set describes a hypothetical household on such a tariff
rather than a real one.

The scenarios follow the factorial design of
Section~\ref{sec:factorial} with one change. The window contains no day
with a negative price, so the set has three price patterns, stable,
volatile and spike, each defined relative to the days of the window.
With two household loads, two peak limits and four days per pattern and
load, the set has 48 scenarios, 14 of which require coordination. The
cost of starting every appliance when it is ready has a median of
\EUR{0.63} per scenario. Every system of the confirmatory evaluation was
run on the set in the same way as on Q2: the Multi-Agent Pipeline and
the LLM without tools three times per scenario with GPT-4o-mini, and
EMHASS and the rule-based systems once.

\section{Results}

Table~\ref{tab:uk-systems} summarises the results. The ranking of the
systems is the same as on the Finnish benchmarks. The pipeline produced
a valid schedule in all 144 runs and saved 6.69\,\% of the cost without
scheduling, the same as the optimal schedule. EMHASS saved 6.45\,\% and
the smart timer 6.55\,\%. The cheapest-hours automations saved 2.6 to
3.6\,\%, and on the 14 scenarios that require coordination only 7\,\% of
their schedules were valid, against 100\,\% for the pipeline, EMHASS and
the smart timer. The LLM without tools saved 0.18\,\%, and a quarter of
its runs cost more than starting every appliance when ready. The
pipeline's mean latency was 18.5 seconds, and its API cost \$0.0025 per
recommendation.

\begin{table}[h!]
\caption{Robustness set: UK prices and tariff, December 2014 to July 2015 (48 scenarios; saving pooled over all runs).}
\label{tab:uk-systems}
\small
\begin{tabularx}{\textwidth}{>{\raggedright\arraybackslash}Xcccc}
\toprule
System & Valid & Saving & \makecell{Saving,\\strict limit} & \makecell{Saving,\\realistic limit} \\
\midrule
Multi-Agent Pipeline & 100\,\% & 6.69\,\% & 6.4\,\% & 6.9\,\% \\
Smart timer (own heuristic) & 100\,\% & 6.55\,\% & 6.2\,\% & 6.9\,\% \\
EMHASS, measured profiles & 100\,\% & 6.45\,\% & 6.0\,\% & 6.9\,\% \\
EMHASS, nominal power & 100\,\% & 6.06\,\% & 5.6\,\% & 6.5\,\% \\
Cheapest window & 70.8\,\% & 3.55\,\% & 1.5\,\% & 5.6\,\% \\
Cheapest-hours automation & 68.8\,\% & 2.62\,\% & 0.9\,\% & 4.3\,\% \\
LLM without tools & 95.1\,\% & 0.18\,\% & 0.3\,\% & 0.1\,\% \\
\midrule
Optimum (upper bound) & 100\,\% & 6.69\,\% & 6.4\,\% & 6.9\,\% \\
\bottomrule
\end{tabularx}
\end{table}

Paired per scenario, the pipeline saved more than the cheapest-hours
automation (by 3.3 percentage points), the cheapest window (2.1) and the
LLM without tools (6.4), with Wilcoxon $p < 0.001$ in each case, and it
saved the same as EMHASS within $\pm 1$ percentage point (mean difference
$+0.20$ points, 90\,\% confidence interval $[0.06, 0.34]$). These tests
mirror the registered ones but were not registered for this set, so they
are reported without adjustment.

The savings are about a quarter of those on the Finnish benchmarks.
Because the price patterns are defined relative to the days of the
window, even the volatile and spike days of 2015 are mild compared with
the Finnish ones: on these days the spot price varied within the day by
8 to 9 cents per kWh on average, against 16 to 20 cents on the
corresponding days of Q2. The window also has no negative prices. In addition,
the constant charge makes up most of the consumer price, so a given
shift in time changes a smaller share of the cost. The order of the
price patterns also differs from the Finnish sets: the stable days gave
the largest saving (9.8\,\% at the optimum) and the volatile days the
smallest (3.0\,\%). The reason was not analysed further.

\section{Conclusion}

Priced at the prices and tariff of their own country and period, the
REFIT appliance cycles give much smaller savings, but the comparison
between the systems does not change. The pipeline again reaches the
optimal saving with a valid schedule every time, matches EMHASS, and
saves more than the cheapest-hours automations, whose schedules again
fail where the household's appliances must be staggered. The main
results of this thesis therefore do not depend on combining UK
appliance data with Finnish prices, although the size of the saving
does depend on the price conditions.


% ════════════════════════════════════════════════════════════════════════════
\chapter{Evaluation of the earlier design}
\label{app:earlier}

This appendix reports the evaluation of the earlier design of the study,
summarised in Section~\ref{sec:earlier-design}. It compared the base
multi-agent pipeline, without retrieval or the optimiser tool, with a
single ReAct agent that had the first seven tools of
Table~\ref{tab:tools} and the same step budget, on the first benchmark of
48 scenarios in six types (Section~\ref{sec:dataset}) and on the
held-out Q set (Section~\ref{sec:resolution}). Its outcome was the
optimality gap, the cost of a schedule above the optimal cost in percent
of the optimal cost, with an infeasible schedule scored at the cost of
starting every appliance when ready. The research questions and
hypotheses of this design were replaced in the redesign, so the labels
SQ1, SQ2, H1 and H2 in Section~\ref{sec:confirmatory} refer to the first
pre-registration. The results remain valid for that design.

\section{Pipeline and single agent on the first benchmark}
\label{sec:statistics}

% Hand-typed from `hems analyze` (results/tables/equivalence.csv, efficiency_tests.csv), GPT-4o-mini only.
This section reports the tests of the earlier design for GPT-4o-mini on
the first benchmark, with the mean over three runs per scenario as the
paired observation ($N = 48$). These results are exploratory; the
pre-registered confirmatory results are reported in
Section~\ref{sec:confirmatory}. A run with the local model gpt-oss-20b
was stopped before completion and is not reported.

\paragraph{Schedule quality.} The multi-agent pipeline reached a mean
optimality gap of 1.25\,\% and the single agent 1.43\,\%. The Wilcoxon
signed-rank test finds no significant difference ($W = 449$,
$Z = -0.14$, $p = 0.89$, $r = 0.02$), and the median paired difference
is zero, because both designs reach the same cost in most scenarios. The
equivalence test shows that this is not merely a lack of evidence. The
mean paired difference is $-0.18$ percentage points, with a 90\,\%
confidence interval of $[-0.71, 0.35]$, which lies inside the primary
margin of $\pm 1$ percentage point (Table~\ref{tab:equivalence}). The two
designs are therefore equivalent in schedule quality within one
percentage point of the optimum. At the narrower margin of 0.5
percentage points the $t$ version of the test does not conclude
equivalence, because a few scenarios with large differences in either
direction widen the interval, although the Wilcoxon version does.

\begin{table}[h!]
\caption{Equivalence test (TOST) of the optimality gap, multi-agent vs.\ single-agent (GPT-4o-mini, $N = 48$).}
\label{tab:equivalence}
\small
\begin{tabular*}{\textwidth}{@{\extracolsep{\fill}}lccc}
\toprule
Margin & $p$, paired $t$ TOST & $p$, Wilcoxon TOST & Equivalent \\
\midrule
$\pm 0.5$ percentage points & 0.156 & $<$\,0.001 & No \\
$\pm 1$ percentage point (primary) & 0.006 & $<$\,0.001 & Yes \\
$\pm 2$ percentage points & $<$\,0.001 & $<$\,0.001 & Yes \\
\bottomrule
\end{tabular*}
\end{table}

\paragraph{Efficiency.} The two designs differ significantly in every
efficiency measure (Table~\ref{tab:efficiency-tests}). The pipeline used
34\,\% fewer input tokens, because each specialist works with a short
context and only its own tools, while the single agent carries the
entire conversation and all eight tool descriptions through every step.
It produced 83\,\% more output tokens and made 23\,\% more tool calls,
since four agents each write their own structured output and call their
tools separately. Input tokens dominate the price, so the API cost per run was
21\,\% lower for the pipeline, and lower in 35 of the 48 scenarios. The
sequential hand-offs made the pipeline 69\,\% slower, with a mean of
21.4 seconds against 12.7 seconds, and it was faster in only two
scenarios.

\begin{table}[h!]
\caption{Efficiency per run, multi-agent vs.\ single-agent (GPT-4o-mini; tokens and tool calls per run, API cost in USD, latency in seconds; each value is the mean over $N = 48$ scenarios of the per-scenario mean of three runs; Wilcoxon signed-rank test).}
\label{tab:efficiency-tests}
\small
\begin{tabular*}{\textwidth}{@{\extracolsep{\fill}}lcccccc}
\toprule
Measure & \makecell{Single\\agent} & \makecell{Multi-\\agent} & Difference & \makecell{Multi-agent\\lower} & $p$ & $r$ \\
\midrule
Input tokens & 21\,686 & 14\,292 & $-34$\,\% & 43/48 & $<$\,0.001 & 0.79 \\
Output tokens & 661 & 1\,208 & $+83$\,\% & 1/48 & $<$\,0.001 & 0.87 \\
Total tokens & 22\,347 & 15\,500 & $-31$\,\% & 40/48 & $<$\,0.001 & 0.75 \\
API cost (USD) & 0.0037 & 0.0029 & $-21$\,\% & 35/48 & $<$\,0.001 & 0.54 \\
Latency (s) & 12.7 & 21.4 & $+69$\,\% & 2/48 & $<$\,0.001 & 0.86 \\
Tool calls & 16.5 & 20.3 & $+23$\,\% & 1/48 & $<$\,0.001 & 0.81 \\
\bottomrule
\end{tabular*}
\end{table}

Taken together, decomposing the task into specialist agents did not
lower the cost of the schedules, but it reached the same quality with
about a third fewer tokens and a fifth lower inference cost, at the price
of a longer response time.

\section{Comparison with baselines and real tools}
\label{sec:baselines}

% Hand-typed from `hems analyze --runs results/runs.jsonl results/baselines.jsonl` (GPT-4o-mini).
Table~\ref{tab:baselines} compares the two agent designs with the smart
timer, the LLM without tools and the two real home energy management
tools. All systems that use a language model use GPT-4o-mini.

% EMHASS rows: Docker/REST service with default settings (tools/emhass_rest.py, systems emhass_rest*); cheapest hours;
% all in results/tools/scenarios.jsonl (secondary, added after the pre-registration). Savings captured = pooled.
\begin{table}[h!]
\caption{Agent designs, simpler baselines and real home energy management tools (GPT-4o-mini, main benchmark, 48 scenarios).}
\label{tab:baselines}
\small
\begin{tabular*}{\textwidth}{@{\extracolsep{\fill}}lcccccc}
\toprule
System & \makecell{Feasible\\runs} & \makecell{Mean\\gap} & \makecell{Median\\gap} & \makecell{Gap,\\type D} & \makecell{Savings\\captured} & \makecell{Time\\per run} \\
\midrule
Smart timer (own heuristic) & 48/48 & 0.23\,\% & 0.00\,\% & 1.2\,\% & 99\,\% & $<$\,0.01\,s \\
EMHASS, profiles & 48/48 & 0.32\,\% & 0.00\,\% & 0.1\,\% & 99\,\% & 1.2\,s \\
Multi-agent pipeline & 144/144 & 1.25\,\% & 0.01\,\% & 5.9\,\% & 93\,\% & 21.4\,s \\
Single agent & 143/144 & 1.43\,\% & 0.00\,\% & 6.5\,\% & 90\,\% & 12.7\,s \\
EMHASS, nominal power & 48/48 & 2.24\,\% & 0.45\,\% & 8.5\,\% & 87\,\% & 1.0\,s \\
Cheapest hours & 25/48 & 17.6\,\% & 4.2\,\% & 62.7\,\% & 25\,\% & $<$\,0.01\,s \\
LLM without tools & 126/144 & 23.1\,\% & 11.9\,\% & 66.0\,\% & 13\,\% & 2.0\,s \\
\bottomrule
\end{tabular*}
\end{table}

Without tools, GPT-4o-mini schedules poorly. Its mean optimality gap is
23.1\,\%, and 18 of its 144 schedules are infeasible: 14 exceed the peak
limit and 4 start a cycle outside its window. A further 22 of its
feasible schedules cost more than starting every appliance as soon as it
is ready, and over all its runs it captures only 13\,\% of the
achievable saving, against about 90\,\% for the agents. Both agent designs are significantly better (Wilcoxon
$p < 0.001$, $r = 0.87$ for each), with median paired differences of
10.8 percentage points for the single agent and 11.4 for the pipeline.
The same model that fails to schedule from a single prompt comes close to
the optimum when it can query prices, costs and loads through tools and
check its own work. Tool use, rather than the language model's own
calculation, is therefore what makes the agents viable.

The smart timer, however, outperforms both agent designs. It finds a
feasible schedule in every scenario, is optimal in 30 of the 48, and its
largest gap is 3.7\,\% (D04). Both agent designs have significantly
larger gaps than the timer (single agent $p < 0.001$, $r = 0.63$;
pipeline $p < 0.001$, $r = 0.68$). The difference is concentrated in
type D, where the timer's mean gap is 1.2\,\% against 5.9\,\% for the
pipeline and 6.5\,\% for the single agent. The median paired differences
are small (0.002 percentage points for the single agent and 0.017 for the
pipeline), because in most scenarios all three systems reach the optimum
or come very close to it. The timer needs no language model, runs in
about two milliseconds and has no inference cost. In this benchmark,
where the power profile and window of every appliance are known exactly
and the horizon is one day, placing the appliances one at a time is
almost always enough. The agents come within about one percentage point
of this rule but do not improve on it. What they add is the ability to
start from a request written in natural language rather than from
structured appliance data, and to explain the schedule they produce.

The two real home energy management tools, added after the
pre-registration, bracket the agents. Configured with the measured power
profiles, EMHASS comes close to the optimum: every schedule is feasible,
32 of the 48 are optimal, and the mean gap is 0.32\,\%. Its largest gap,
4.2\,\% in F02, comes from its default solver tolerance of 1\,\%
(the setting \texttt{lp\_solver\_mip\_rel\_gap} = 0.01; EMHASS
contributors, 2026), which
applies to the whole electricity bill including the base load and can
therefore amount to several percent of the appliance cost scored here.
With an exact tolerance the same optimization comes within 0.2\,\% of
the optimum in every scenario, which confirms that the MILP reference reflects what a
deployed open-source optimizer can achieve. The agents' mean gaps are
about one percentage point higher, but the difference from EMHASS as
deployed is not significant (Wilcoxon $p = 0.24$ for the single agent and
$p = 0.19$ for the pipeline). Configured with a constant nominal power,
EMHASS still keeps every schedule within the peak limit, because it
plans each cycle at its peak power, but it misjudges the cost of the real
cycles and reaches a mean gap of 2.2\,\%. Both agent designs are
significantly better than this configuration (Wilcoxon $p < 0.001$ for
each, $r = 0.67$ for the single agent and $0.64$ for the pipeline). The
cheapest-hours automation, under the assumption stated in
Section~\ref{sec:evaluation} that its trigger starts a cycle which then
runs to completion, exceeds the peak limit in 23 of the 48 scenarios,
because each appliance follows
its own cheap period and appliances whose windows share a price minimum
start together. Scored at the naive cost when infeasible, it reaches a
mean gap of 17.6\,\%, and both agent designs are significantly better
($p < 0.001$, $r = 0.82$ and $0.81$). The agents are therefore
outperformed by the smart timer and are not significantly different from
a well-configured optimizer as deployed, while they outperform a simply
configured optimizer and the per-device automation that is common among
households.

Two checks show that this ranking does not depend on how the costs are
scored. First, scoring an infeasible run at the cost of the naive
schedule does not drive the result. Restricted to feasible runs, the
mean gaps are 0.23\,\% for the smart timer, 0.32\,\% for EMHASS with
profiles, 1.25\,\% for the pipeline, 1.45\,\% for the single agent,
2.24\,\% for EMHASS with nominal power, 7.9\,\% for the cheapest-hours
automation and 23.0\,\% for the LLM without tools, the same order as in
Table~\ref{tab:baselines}, and the order on the held-out benchmark is
also unchanged. The cheapest-hours automation looks better on feasible
runs only because these are the days on which its appliances happen not
to collide. Second, the ranking is the same under the spot price alone
as under the consumer price, as explained in
Section~\ref{sec:evaluation}.

\section{Confirmatory run on the held-out Q set}
\label{sec:confirmatory}

% Hand-typed from `hems confirm` (results/confirmation/confirmatory_tests.csv) and
% `hems analyze --runs results/confirmation.jsonl` for the secondary results.
The hypotheses of the earlier SQ1 and SQ2 were fixed in the pre-registration, which
was committed to the project repository on 30 September 2026 together
with the code, prompts, tools and analysis script. The confirmatory run
followed on the held-out benchmark of 48 scenarios with fifteen-minute
prices (Section~\ref{sec:resolution}), on which no language model had
been run before. All systems that use a language model used GPT-4o-mini
with three runs per scenario, and the analysis was run exactly as
registered. The run cost USD~0.92 in API fees. Table~\ref{tab:confirmatory}
reports the six registered tests.

\begin{table}[h!]
\caption{Pre-registered confirmatory tests on the held-out fifteen-minute benchmark (GPT-4o-mini, $N = 48$; $p$-values Holm-adjusted within each sub-question).}
\label{tab:confirmatory}
\small
\begin{tabularx}{\textwidth}{l>{\raggedright\arraybackslash}Xcccc}
\toprule
& Hypothesis & Estimate & $p$ & $r$ & Supported \\
\midrule
H1 & Single agent gap $<$ LLM without tools & $-17.9$ pp & $<$\,0.001 & 0.86 & Yes \\
H1 & Pipeline gap $<$ LLM without tools & $-17.7$ pp & $<$\,0.001 & 0.86 & Yes \\
H2a & Pipeline equivalent to single agent ($\pm 1$ pp) & $+0.22$ pp & $<$\,0.001 & -- & Yes \\
H2b & Pipeline uses fewer input tokens & $-31$\,\% & $<$\,0.001 & 0.84 & Yes \\
H2b & Pipeline has a lower API cost per run & $-21$\,\% & $<$\,0.001 & 0.70 & Yes \\
H2c & Pipeline has a longer latency & $+52$\,\% & $<$\,0.001 & 0.80 & Yes \\
\bottomrule
\end{tabularx}
\end{table}

All six hypotheses are supported. For H1, both agent designs reached a
mean optimality gap of about one percent (0.86\,\% for the single agent
and 1.08\,\% for the pipeline), whereas the LLM without tools reached
18.8\,\%. Fifteen of its 144 schedules were infeasible, and over all its
runs it captured only 24\,\% of the achievable saving, against 92\,\% for
the single agent and 90\,\% for the pipeline. For H2a,
the mean paired difference in gap was $+0.22$ percentage points with a
90\,\% confidence interval of $[-0.15, 0.59]$, which lies well inside
the margin of $\pm 1$ percentage point. On this benchmark the pipeline
was thus marginally worse than the single agent on average, the reverse
of the exploratory result, but in both runs the difference is small
enough to count as equivalent. For H2b and H2c, the pipeline used
13\,310 input tokens per run against 19\,189 for the single agent, its API
cost was 21\,\% lower, and its latency was 20.4 seconds against 13.4
seconds.

The confirmatory estimates are close to the exploratory ones:
input tokens fell by 31\,\% in both runs, API cost by 21\,\% in both, and
latency rose by 52\,\% here against 69\,\% before. The secondary results
also replicate. The equivalence test is again not significant at the
narrower margin of $\pm 0.5$ percentage points in its $t$ version
($p = 0.11$), although its Wilcoxon version is. The smart timer again
outperforms both agent designs, with a mean gap of 0.17\,\% (Wilcoxon
$p < 0.001$ against each design, $r = 0.69$ for the single agent and
$0.64$ for the pipeline). Type D remains the hardest scenario type, with
mean gaps of 3.9\,\% for the single agent and 5.5\,\% for the pipeline,
while the LLM without tools fails most clearly in type E, with a mean gap
of 56\,\%.

The real home energy management tools were added after the
pre-registration, so the following comparisons are secondary and not
registered; their $p$-values are descriptive and not adjusted for
multiple testing. EMHASS as deployed reaches a mean gap of 0.36\,\% with
measured profiles, not significantly different from either agent design
(Wilcoxon $p = 0.90$ for the single agent and $p = 0.69$ for the
pipeline), and 3.5\,\% with nominal power, significantly worse than both
($p < 0.001$). The cheapest-hours automation exceeds the peak limit in 26
of the 48 scenarios, with a mean gap of 18.7\,\%. These patterns match
those on the main benchmark.

The registered findings, H1 for SQ1 and H2 for SQ2, are therefore not
specific to the benchmark on which the pipeline was developed: they hold
on unseen days with a different price resolution. They cover GPT-4o-mini
only. The registered hypothesis on a small model run locally (H3) was
withdrawn in the first amendment of the pre-registration, before its run
on the first benchmark was completed.

% ── References cited only in Appendix A: move back into the References list ──
% (they are inside \begin{hangparas} ... \end{hangparas} in main.tex)
%
% Department of Energy \& Climate Change. (2016). Domestic energy bills in
% 2015: The impact of variable consumption. In \textit{Energy Trends:
% March 2016}. UK Government. Retrieved 6 October 2026, from
% \url{https://assets.publishing.service.gov.uk/media/5a7f857740f0b6230268ffdc/Domestic_energy_bills_in_2015_-_The_impact_of_variable_consumption.pdf}
% 
% European Central Bank. (2026). \textit{Euro foreign exchange reference
% rates: Pound sterling (GBP)} [Data set]. Retrieved 6 October 2026, from
% \url{https://data.ecb.europa.eu/data/datasets/EXR}
